% !PRINTED 157
\weberpage{157}

%<<CONT
in bezug auf $u$ integrieren, so folgt:

\weq{1}{\dfrac{1}{2}\log\dfrac{\Theta(u-v)}{\Theta(u+v)}+u\,Z(v) = \int_0^u \dfrac{\varkappa^2\,\mathrm{sn}^2u\,\mathrm{sn}\,v\,\mathrm{cn}\,v\,\mathrm{dn}\,v}{1-\varkappa^2\,\mathrm{sn}^2u\,\mathrm{sn}^2v}\,du,}

und wir können noch die auf gleiche Weise [aus § 43 (11)]
herzuleitende Formel:

\weq{2}{\dfrac{1}{2}\log\dfrac{H(v-u)}{H(v+u)}+u\,Z(v) = \int_0^u \dfrac{\mathrm{sn}\,v\,\mathrm{cn}\,v\,\mathrm{dn}\,v\,du}{\mathrm{sn}^2u-\mathrm{sn}^2v}} % Source recheck: damaged upper-limit sort retained as u; du and the parallel formulas (1), (3) support this reading.

hinzufügen, die übrigens auch aus (1) abgeleitet werden kann.
Wir setzen nun mit \gesperrt{Jacobi}

\weq{3}{\Pi(u,v) = \int_0^u \dfrac{\varkappa^2\,\mathrm{sn}^2u\,\mathrm{sn}\,v\,\mathrm{cn}\,v\,\mathrm{dn}\,v}{1-\varkappa^2\,\mathrm{sn}^2u\,\mathrm{sn}^2v}\,du,}

und nennen $\Pi(u,v)$ \gesperrt{die Transzendente dritter
Gattung} mit dem Argument $u$ und dem Parameter $v$. Ersetzt
man $du$ durch den Ausdruck:

\[
\dfrac{d\xi}{2\sqrt{\xi(1-\xi)(1-\varkappa^2\xi)}}
\]

und $\mathrm{sn}^2u$ durch $\xi$, so ergibt sowohl (1) als (2)
ein elliptisches Integral dritter Gattung, wie wir es in § 11
kennen gelernt haben.

Aus (1) folgt zunächst

\weq{4}{\Pi(u,v)-u\,Z(v) = \Pi(v,u)-v\,Z(u),}

oder der \gesperrt{Jacobische Satz} über die \gesperrt{Vertauschung von
Argument und Parameter}.

Die Funktion $\Pi(u,v)$ verschwindet, wenn $v=K$ oder
$v=K+iK'$ ist, weil für den ersteren Wert $\mathrm{cn}\,v$, für
den zweiten $\mathrm{dn}\,v$ gleich Null ist. Setzen wir daher
diese Werte für $u$ in (4) ein, so folgt

\weq{5}{\begin{aligned}
\Pi(K,v) &= K\,Z(v)\\
\Pi(K+iK',v) &= (K+iK')\,Z(v)+\dfrac{i\pi v}{2K}\ [\text{§ 47, (13)}],
\end{aligned}}

wodurch die vollständigen Integrale dritter Gattung auf die
zweite Gattung zurückgeführt sind.

Wir wollen noch das Additionstheorem der $\Pi$-Funktion
ableiten, das \gesperrt{Jacobi} in den Fund. nova art. 53--55\footnote{C.\,G.\,J.\ \gesperrt{Jacobis} gesammelte Werke, Berlin 1881, Bd.\,I, S.\,204.} in ver-
%CONT>>
